\section{Experiments}
\label{sec:experiments}

This section describes experimental setup, diagnostics, and preliminary explorations that motivated the final controlled comparison.

\subsection{Experiment E1a: Native Video Diagnostic}

\textbf{Objective}: Verify that the base model receives true temporal input and establish performance ceiling under native video processing.

\textbf{Setup}:
\begin{itemize}
    \item Model: Qwen2.5-VL-3B-Instruct (base, no fine-tuning)
    \item Input: Native Qwen video pathway (not separate-image fallback)
    \item Frames: Exactly 24 per video
    \item All 50 card-relocation videos
    \item Verification: Check that native temporal metadata (second\_per\_grid\_ts) present for all records
\end{itemize}

\textbf{Results}:
\begin{itemize}
    \item Valid JSON: 50/50 records
    \item Complete six-round trace (exact match): 0/50 (0\%)
    \item Final-position accuracy: 12/50 (24\%)
    \item Single-round position accuracy: \textasciitilde19\%
    \item Native temporal metadata verified: all 50 records have 24 frames
\end{itemize}

\textbf{Interpretation}: Base model receives true temporal input but achieves near-random accuracy on final position (24\% vs. 25\% random baseline) and 0\% on complete traces. This establishes that poor performance is not due to input processing failures---the model has genuine difficulty tracking state.

\subsection{Baseline (Zero-Shot) Evaluation on Test Set}

\textbf{Results on 10 held-out videos}:
\begin{itemize}
    \item Final-answer prompt: 0.20 (2/10)
    \item Trajectory prompt: 0.20 (2/10)
    \item Per-round accuracy: 0.167 (10/60)
\end{itemize}

Baseline performance establishes a floor: any improvement from supervised training should exceed this.

\subsection{Final-Frame-Only Baseline}

\textbf{Objective}: Control for static image shortcuts (e.g., final frame revealing position).

\textbf{Setup}: Run model on final frame only (without prior video context).

\textbf{Results on 50 videos}:
\begin{itemize}
    \item Accuracy: 12/50 (24\%)
    \item Random baseline: 25\%
    \item Majority baseline: 30\%
    \item Prediction distribution: A=0, B=32, C=12, D=6
    \item Ground-truth distribution: A=11, B=14, C=10, D=15
\end{itemize}

Conclusion: Final-frame-only does not exceed baselines. Models cannot exploit static image shortcuts.

\subsection{Diagnostic Findings}

\textbf{Discovery: target\_card $\equiv$ initial\_position}

Audit across all 50 videos revealed 100\% of videos have target card in initial position. Early initial-position diagnostics achieved \textasciitilde76\% accuracy, but this was spurious: models could simply emit the target card label without visual grounding. This confound invalidated early diagnostics and motivated redesigned evaluation (final-position, trajectory, exact-match metrics).

\textbf{Invalid L2 Experiment}

An earlier experiment (``L2'') used ground-truth files with incorrect field mapping (position\_after\_round\_X instead of position\_trace). Results: 0/50 accuracy. This experiment was rejected from scientific results and documented in validity audit section.

\subsection{Exploratory Frame-Order Control (Step 9)}

\textbf{Objective}: Exploratory test whether temporal ordering matters.

\textbf{Setup}: Model B evaluated on shuffled video frames.

\textbf{Important limitation}: This experiment used \emph{separate-image input}, not native video. Technical constraint discovered post-hoc. Therefore, results cannot isolate temporal-order effects cleanly.

\textbf{Results}:
\begin{itemize}
    \item True order final accuracy: 0.10 (1/10)
    \item Shuffled order final accuracy: 0.10 (1/10)
    \item Per-round: 0.25 vs. 0.25 (identical)
\end{itemize}

\textbf{Status}: Retained for transparency, marked as exploratory and unsuitable for causal claims about temporal ordering. Recommendation for future work: implement frame shuffling within native video pipeline for proper testing.
